\SixthArticle{27}
The integral $\int x^{\lambda-1}(1-x^\mu)^\nu\,\mathrm{d}x$, chosen so that it vanishes for $x=0$, is expressed by the following series, provided $\lambda,\mu$ are positive quantities:
\[
\frac{x^\lambda}{\lambda}
-\frac{\nu x^{\mu+\lambda}}{\mu+\lambda}
+\frac{\nu(\nu-1)x^{2\mu+\lambda}}{1.2.(2\mu+\lambda)}
-\mathrm{etc.}
=\frac{x^\lambda}{\lambda}F(-\nu,\frac\lambda\mu,
\frac\lambda\mu+1,x^\mu)
\]
Hence its value for $x=1$ will be
\[
=\frac{\Pi\frac\lambda\mu.\Pi\nu}
{\lambda\Pi(\frac\lambda\mu+\nu)}.
\]
All the relations that the distinguished \textsc{Euler} once worked out with much labour follow of themselves from this theorem. Thus, for example, on setting
\[
\int\frac{\mathrm{d}x}{\BareRoot{(1-x^4)}}=A,\qquad
\int\frac{xx\,\mathrm{d}x}{\BareRoot{(1-x^4)}}=B
\]
the result is $A=\frac{\Pi\frac14.\Pi(-\frac12)}{\Pi(-\frac14)}$, $B=\frac{\Pi\frac34.\Pi(-\frac12)}{3\Pi\frac14}
=\frac{\Pi(-\frac14).\Pi(-\frac12)}{4\Pi\frac14}$, and therefore $AB=\frac14\pi$. At the same time it follows, since $\Pi\frac14.\Pi(-\frac14)=\frac{\frac14\pi}{\sin\frac14\pi}
=\frac\pi{\BareRoot{8}}$, that
\[
\Pi\frac14=\BareRoot[4]{(\frac18\pi AA)}
=\BareRoot[4]{\frac{\pi^3}{128BB}},\qquad
\Pi(-\frac14)=\BareRoot[4]{\frac{\pi^3}{8AA}}
=\BareRoot[4]{(2\pi BB)}
\]
The numerical value of $A$, as calculated by \textsc{Stirling}, is $=1{,}3110287771\;\allowbreak 4605987$; the value of $B$, according to the same author, is $=0{,}5990701173\;\allowbreak 6779611$, and according to our calculation, based on a special device, $=0{,}5990701173\;\allowbreak 6779610372$.

In general, it is easy to show that the value of the function $\Pi z$, if $z$ is a rational quantity $=\frac m\mu$, with $m,\mu$ denoting integers, can be deduced from $\mu-1$ determined values of such integrals for $x=1$, and indeed in very many different ways. For if an integer is taken for $\lambda$ and for $\nu$ a fraction whose denominator is $=\mu$, the value of that integral always reduces to three $\Pi z$, where $z$ is a fraction with denominator $=\mu$; any such $\Pi z$, however, can be reduced by formula 45 either to $\Pi(-\frac1\mu)$, or to $\Pi(-\frac2\mu)$, or to $\Pi(-\frac3\mu)$ etc., or to $\Pi(-\frac{\mu-1}{\mu})$, provided $z$ is actually a fraction; for if $z$ is an integer, $\Pi z$ is known in itself. From those values of the integrals, however, generally speaking, any $\Pi(-\frac m\mu)$,
\SourceLeaf{160}
if $m<\mu$, can be obtained by elimination\HistoricalNote{17}.
Indeed, even half as many such integrals will suffice if we also call on formula 54 for help. Thus, for example, on setting
\[
\begin{gathered}
\int\frac{\mathrm{d}x}{\BareRoot[5]{(1-x^5)}}=C,\quad
\int\frac{\mathrm{d}x}{\BareRoot[5]{(1-x^5)^2}}=D,\quad
\\\int\frac{\mathrm{d}x}{\BareRoot[5]{(1-x^5)^3}}=E,\quad
\int\frac{\mathrm{d}x}{\BareRoot[5]{(1-x^5)^4}}=F,
\end{gathered}
\]
the result is
\[
\begin{gathered}
C=\Pi\frac15.\Pi(-\frac15),\quad
D=\frac{\Pi\frac15.\Pi(-\frac25)}{\Pi(-\frac15)},\quad
\\E=\frac{\Pi\frac15.\Pi(-\frac35)}{\Pi(-\frac25)},\quad
F=\frac{\Pi\frac15.\Pi(-\frac45)}{\Pi(-\frac35)}
\end{gathered}
\]
Hence, since $\Pi\frac15=\frac15\Pi(-\frac45)$, we have
\[
\begin{gathered}
\Pi(-\frac15)=\BareRoot[5]{\frac{5C^4}{DEF}},\quad
\Pi(-\frac25)=\BareRoot[5]{\frac{25C^3D^3}{EEFF}},\quad
\\\Pi(-\frac35)=\BareRoot[5]{\frac{125CCDDEE}{F^3}},\\
\Pi(-\frac45)=\BareRoot[5]{(625CDEF)}
\end{gathered}
\]
Formulas 54 and 55 further give
\[
C=\frac\pi{\sin\frac15\pi},\qquad
\frac DF=\frac{\sin\frac15\pi}{\sin\frac25\pi}
\]
so that the two integrals $D,E$, or $E$ and $F$, suffice to calculate all the values $\Pi(-\frac15)$, $\Pi(-\frac25)$ etc.

\SixthArticle{28}
On setting $y=\nu x$ and $\mu=1$, $\frac{\Pi\lambda.\Pi\nu}{\lambda\Pi(\lambda+\nu)}$ will be the value of the integral $\int\frac{y^{\lambda-1}(1-\frac y\nu)^\nu\,\mathrm{d}y}{\nu^\lambda}$ from $y=0$ to $y=\nu$, or the value of the integral $\int y^{\lambda-1}(1-\frac y\nu)^\nu\,\mathrm{d}y$ between the same limits will be $=\frac{\nu^\lambda\Pi\lambda.\Pi\nu}{\lambda\Pi(\lambda+\nu)}
=\frac{\Pi(\nu,\lambda)}\lambda$ (form. 47), provided $\nu$ denotes an integer. Now, as $\nu$ grows to infinity, the limit of $\Pi(\nu,\lambda)$ will be $=\Pi\lambda$, while the limit of $(1-\frac y\nu)^\nu$ will be $e^{-y}$, with $e$ denoting the base of natural logarithms. Therefore, if $\lambda$ is positive, $\frac{\Pi\lambda}\lambda$ or $\Pi(\lambda-1)$ will express the integral $\int y^{\lambda-1}e^{-y}\,\mathrm{d}y$ from $y=0$ to $y=\infty$; or, writing $\lambda$ for $\lambda-1$, $\Pi\lambda$ is the value of the integral $\int y^\lambda e^{-y}\,\mathrm{d}y$ from $y=0$ to $y=\infty$, if $\lambda+1$ is a positive quantity.

More generally, on setting $y=z^\alpha$, $\alpha\lambda+\alpha-1=\beta$, the integral $\int y^\lambda e^{-y}\,\mathrm{d}y$ becomes $\int\alpha z^\beta e^{-z^\alpha}\,\mathrm{d}z$, which, taken between the limits $z=0$ and $z=\infty$, will therefore be expressed by $\Pi(\frac{\beta+1}\alpha-1)$

\noindent or

The value of the integral $\int z^\beta e^{-z^\alpha}\,\mathrm{d}z$ from $z=0$ to $z=\infty$ is $=\frac{\Pi(\frac{\beta+1}\alpha-1)}\alpha
=\frac{\Pi\frac{\beta+1}\alpha}{\beta+1}$, provided $\alpha$ and $\beta+1$ are positive quantities (if both are negative, the\SourceLeaf{161} integral will be expressed by $-\frac{\Pi\frac{\beta+1}\alpha}{\beta+1}$). Thus, for example, for $\beta=0$, $\alpha=2$, the value of the integral $\int e^{-zz}\,\mathrm{d}z$ is found to be $=\Pi\frac12=\frac12\BareRoot\pi$.

\SixthArticle{29}
For the sum of logarithms $\log1+\log2+\log3+\mathrm{etc.}+\log z$, the distinguished \textsc{Euler} obtained the series $(z+\frac12)\log z-z+\frac12\log2\pi
+\frac{\mathfrak A}{1.2z}-\frac{\mathfrak B}{3.4z^3}
+\frac{\mathfrak C}{5.6z^5}-\mathrm{etc.}$, where $\mathfrak A=\frac16$, $\mathfrak B=\frac1{30}$, $\mathfrak C=\frac1{42}$ etc. are the \textsc{Bernoulli} numbers. Thus $\log\Pi z$ is expressed by this series; for although at first sight this conclusion may seem restricted to integer values, closer consideration shows that the expansion used by \textsc{Euler} (Instit. Calc. Diff. Cap.\,\textsc{vi}. 159) can be applied at least to positive fractional values with the same right as to integers: it assumes only that the function of $z$ to be expanded into a series is such that its decrease, when $z$ passes to $z-1$, can be expressed by \textsc{Taylor}'s theorem, and at the same time that this decrease is $=\log z$. The first condition rests on the \emph{continuity} of the function, and therefore does not hold for negative values of $z$, to which that series must consequently not be extended; the second condition, however, holds for the function $\log\Pi z$ in general, without restriction to integer values of $z$. Let us therefore set
\[
\begin{aligned}
{}\text{[58]}\quad &\log\Pi z=(z+\frac12)\log z-z+\frac12\log2\pi\\[4pt]
&+\frac{\mathfrak A}{1.2z}-\frac{\mathfrak B}{3.4z^3}
+\frac{\mathfrak C}{5.6z^5}-\frac{\mathfrak D}{7.8z^7}
+\mathrm{etc.}
\end{aligned}
\]
Since this also gives
\[
\begin{aligned}
&\log\Pi2z=(2z+\frac12)\log2z-2z+\frac12\log2\pi\\[4pt]
&+\frac{\mathfrak A}{1.2.2z}-\frac{\mathfrak B}{3.4.8z^3}
+\frac{\mathfrak C}{5.6.32z^5}-\frac{\mathfrak D}{7.8.128z^7}
+\mathrm{etc.}
\end{aligned}
\]
and, by formula 57, on setting $n=2$,
\[
\log\Pi(z-\frac12)=\log\Pi2z-\log\Pi z
-(2z+\frac12)\log2+\frac12\log2\pi,\quad\text{the result is}
\]
\[
\begin{aligned}
{}\text{[59]}\quad &\log\Pi(z-\frac12)=z\log z-z+\frac12\log2\pi\\[4pt]
&-\frac{\mathfrak A}{1.2.2z}+\frac{7\mathfrak B}{3.4.8z^3}
-\frac{31\mathfrak C}{5.6.32z^5}+\frac{127\mathfrak D}{7.8.128z^7}
-\mathrm{etc.}
\end{aligned}
\]
For large values of $z$, these two series initially converge quite rapidly, so that an approximate sum can be obtained conveniently and with sufficient accuracy. Yet it must be carefully noted that for any given value of $z$, however large, only a limited precision can be obtained, since the \textsc{Bernoulli} numbers form a hypergeometric series, and therefore those series, if only extended far enough, certainly pass from convergence to divergence. Besides, it cannot be denied that the theory of such divergent series is still burdened by certain difficulties, which we may perhaps discuss at greater length on another occasion.
